\documentclass[a4paper,12pt]{article}

% Set margins
\usepackage[hmargin=2.5cm, vmargin=3cm]{geometry}

\frenchspacing

% Language packages
\usepackage[utf8]{inputenc}
\usepackage[T1]{fontenc}
\usepackage[magyar]{babel}

% AMS
\usepackage{amssymb,amsmath}

% Graphic packages
\usepackage{graphicx}

% Colors
\usepackage{color}
\usepackage[usenames,dvipsnames]{xcolor}

% Enumeration
\usepackage{enumitem}

% Links
\usepackage{hyperref}

\pagestyle{empty}

\begin{document}

\begin{center}
	{\Large ÜTEMTERV -- \texttt{GEMAK211-B2}}

	\bigskip

	{\huge \textbf{Programtervezési ismeretek}}
	
	\medskip
	
	{\large Programtervező informatikus BSc, Gazdaságinformatikus BSc
	
	nappali tagozat, őszi félév}
\end{center}

\vskip 1cm

\noindent \textbf{1. hét}: Információ, adat, Déscartes szorzat, elemi függvények és műveletek, számrendszerek, számrendszeri átváltás

\bigskip

\noindent \textbf{2. hét}: Számábrázolás: előjeles és előjel nélküli számok ábrázolása, lebegőpontos számábrázolás

\bigskip

\noindent \textbf{3. hét}: Logikai adattípus és műveletei, diszjunktív- és konjunktív normálforma, logikai kapuáramkörök

\bigskip

\noindent \textbf{4. hét}: Adatstruktúrák: tömbök, vektorok, halmaz, karakter, sztring, dátum és idő ábrázolása

\bigskip

\noindent \textbf{5. hét}: Algoritmusok lejegyzési módjai, folyamatábra, pszeudó kód, UML szabvány ide vonatkozó elemei

\bigskip

\noindent \textbf{6. hét}: Elemi algoritmusok, iteráció, rekurzió

\bigskip

\noindent \textbf{7. hét}: Strukturált programozás, Böhm-Jacopini tétel, ciklikus bonyolultság

\bigskip

\noindent \textbf{8. hét}: A számítógép elvi felépítése, fordítás és interpretálás, gépi kód, tárgykód

\bigskip

\noindent \textbf{9. hét}: Többdimenziós adatszerkezetek: mátrixok, listák, fák, fájlformátumok 

\bigskip

\noindent \textbf{10. hét}: Procedurák, függvények, értékátadás, láthatóság, $\lambda$-kalkulus

\bigskip

\noindent \textbf{11. hét}: Típus konverziók: szám-szöveg, szöveg-szám, lebegőpontos számok, dátum és idő konverziója

\bigskip

\noindent \textbf{12. hét}: Hibakezelés: hibák típusai, típusbiztonság, explicit- és implicit hibakezelés, tesztek, assertion-ök

\bigskip

\noindent \textbf{13. hét}: Forráskód szervezése, függőségek kezelése, rendszer integráció, verziókezelés, COCOMO modell

\bigskip

\noindent \textbf{14. hét}: Elvek, módszertanok (DRY, SOLID, SCRUM, KANBAN, TDD)

\vskip 1cm

\noindent \textit{Megjegyzés: A zárthelyi dolgozatok megírására előadások idejében kerül majd sor, emiatt a témakörök bemutatásának időpontja eltérhet a tervezettől.}

\newpage

\noindent \textbf{Az aláírás megszerzésének feltétele}

\smallskip

\noindent A félév közben a két zárthelyi dolgozat lesz (terv szerint a 7. és a 12. héten). Az aláírás megszerzéséhez ezen dolgozatok legalább elégséges szintű megírása szükséges. A dolgozatokat a pótlási héten külön-külön lehet pótolni/javítani. (Egyazon alkalommal a kettőt, tehát akinek egyik sem sikerült előtte, annak mindkét dolgozatot meg kell írnia a rendelkezésre álló időn belül.)

\vskip 1cm

\noindent \textbf{Megajánlott jegy}

\smallskip

\noindent A tárgyból van lehetőség megajánlott jegyet szerezni. Akinek a zárthelyi dolgozataira kapott jegyeinek az átlaga legalább 4.5, az a vizsgát kiválthatja egy, a tárgyhoz kapcsolódó témában készített feladattal, annak bemutatásával.
\begin{itemize}
	\item A témakör az illetővel az érdeklődési körének megfelelően egyeztetésre kerül.
	\item Az aláíráspótlás során javítható a korábbi zárthelyi eredménye.
\end{itemize}



\vskip 1cm

\noindent \textbf{A vizsga}

\smallskip

\noindent A vizsga írásbeli lesz, mely egyaránt tartalmaz elméleti és gyakorlati kérdéseket is.

\vskip 1cm

\noindent \textbf{Ponthatárok}

\smallskip

\noindent A dolgozatokon maximálisan 12 pontot lehet szerezni. Az érdemjegyekre az alábbi ponthatárok vonatkoznak.

\begin{center}
\begin{tabular}{|c|c|}
pont & érdemjegy \\
\hline
0-5 & 1 \\
6 & 2 \\
7-8 & 3 \\
9-10 & 4 \\
11-12 & 5 \\
\hline
\end{tabular}
\end{center}

\vskip 18mm

\hskip 10cm Piller Imre

\hskip 7cm Alkalmazott Matematikai Intézeti Tanszék

\vskip 5mm

\noindent Miskolc, 2025. szeptember 4.

\end{document}

